\section{Objectifs}

    \begin{verbatim}
- Objectif naif : Identifier les composants du modèle impliqués dans l'identification et la gestion 
de la langue de spécialité d'un dataset
- Objectif pratique : Identifier et proposer des préconisations d'usage de la librairie MechIR pour 
la localisation des composant impliqués dans la langue de spécialité
- Définition de langue de spécialité, réduction de l'objectif à l'identification du vocabulaire 
de spécialité 
    \end{verbatim}


\section{Méthodologie}
    \subsection{Construction des perturbations}
        \begin{verbatim}
- on s'intéresse aux perturbations de type replace pour travailler sur l'impact des mots importants
(car liés au vocabulaire de spécialité) en les replaçant
- On propose plusieurs candidats au remplacement classé en 4 niveaux ...
- Le choix des remplaçant est fait par jugement humain
- Le choix des mots à remplacer se fait à l'aide de 
    - fréquence dans l'échantillon de données sélectionné
    - IDF 
    - appartenance au vocabulaire de spécialité (jugement humain)
    - présence du mot dans la question étudiée (cas du dataset à 1! question)
- pour augmenter le nombre de document perturbé et limiter l'impact d'un mot seulement dans la 
cartographie, on propose un nouveau type de perturbation : full_replace
        \end{verbatim}

    \subsection{Estimation de la pertinence d'une perturbation}
        \begin{verbatim}
SI PAS DÉJÀ PRÉSENTÉ DANS PARTIE DE L'ÉTAT DE L'ART SUR MECHIR
- On visualise des scores de similarité entre questoin et document
- graphique d'exemple + texte d'accompagnement à la lecture
- Un perturbation peut être considérée comme pertinente lorsqu'elle modifie significativement 
(jugement humain) la distribution des cores pour l'échantillon du dataset étudié
        \end{verbatim}

    \subsection{Interprétation des résultats de l'Activation Patching}

        \begin{verbatim}
SI PAS DÉJÀ FAITE DANS SOTA sur MechIR
- graphique + accompagnement à la lecture (on obtient une sensibilité de chaque 
noeuds pour chaque paire (Question, document), on visualise la moyenne)

- étant donné la construction de ces valeurs, on s'intéresse aussi aux écarts-types
- pour faciliter la visualisation et la comparaison en cas de sensibilité variable 
d'une perturbation à l'autre on propose de rescaler les poids des matrices
- L'identification des noeuds sensibles doit reposer sur un critère commun. On propose
un critère de sensibilité ...
- Visualisation des noeuds sensibles commun à plusieurs perturbation (matrice des noeuds sensibles)
    -> croiser les différentes caractéristiques (mot remplacé, type, niveau)
    -> critique des noeuds sensibles de la dernière couche
        \end{verbatim}

\section{Résultats et Analyse}

    \subsection{Application à la cartographie du vocabulaire de la défense français dans mE5}
        \begin{verbatim}
- Grand nombre de noeuds de la dernière couche du modèle
- pas de noeuds clairement associé au vocabulaire
...
        \end{verbatim}


    \subsection{Préconisations d'utilisation}
        \begin{verbatim}
- des perturbations qui couvrent tous ou très grande partie des documents étudiés
- utiliser la question comme document témoin
        \end{verbatim}


\section{Perspectives d'évolution de la méthode proposée}

    \begin{verbatim}
- échelle pour situer le niveau d'un mot candidat remplaçant en fonction du mot remplacé
- critère d'identification du vocabulaire de spécialité
- intégration d'autres informations au critère de sensibilité
    \end{verbatim}